
\documentclass{article}
\usepackage{colortbl}
\usepackage{makecell}
\usepackage{multirow}
\usepackage{supertabular}

\begin{document}

\newcounter{utterance}

\centering \large Interaction Transcript for game `codenames', experiment `unambiguous', episode 1 with reflexion{-}Qwen3.5{-}27B.
\vspace{24pt}

{ \footnotesize  \setcounter{utterance}{1}
\setlength{\tabcolsep}{0pt}
\begin{supertabular}{c@{$\;$}|p{.15\linewidth}@{}p{.15\linewidth}p{.15\linewidth}p{.15\linewidth}p{.15\linewidth}p{.15\linewidth}}
    \# & $\;$A & \multicolumn{4}{c}{Game Master} & $\;\:$B\\
    \hline

    \theutterance \stepcounter{utterance}  
    & & & \multicolumn{2}{p{0.3\linewidth}}{
        \cellcolor[rgb]{0.9,0.9,0.9}{
            \makecell[{{p{\linewidth}}}]{
                \texttt{\tiny{[GM$|$GM]}}
                {'hidden': {'team': ['geisha', 'spinach', 'cactus'], 'innocent': ['puck', 'cavalry', 'bayou', 'courtroom', 'luncheon'], 'opponent': [], 'assassin': ['bridegroom']}, 'revealed': {'team': {'team': [], 'innocent': [], 'opponent': [], 'assassin': []}, 'opponent': {'team': [], 'innocent': [], 'opponent': [], 'assassin': []}}}
            }
        }
    }
    & & \\ \\

    \theutterance \stepcounter{utterance}  
    & & \multicolumn{4}{p{0.6\linewidth}}{
        \cellcolor[rgb]{0.9,0.9,0.9}{
            \makecell[{{p{\linewidth}}}]{
                \texttt{\tiny{[P1$\langle$GM]}}
                \texttt{You have attempted this same task before and did not succeed. These are the plans you wrote after those attempts.} \\
\texttt{Trial \#1: New plan: My previous strategy failed because I prioritized grouping multiple words per clue (e.g., "plant" for spinach and cactus) over safety, which led to my partner guessing opponent or distractor words and ending my turn prematurely. I also failed to account for the high risk of ambiguous clues like "view" or "fake" triggering incorrect guesses. Next time, I will adopt a conservative "one{-}word{-}at{-}a{-}time" strategy for the remaining turns. I will select clues that are highly specific to a single remaining target word to minimize the chance of my partner guessing an opponent's word or a distractor. For example, instead of "plant" for spinach and cactus, I would have used "leafy" specifically for spinach or "spine" specifically for cactus, ensuring the clue is too narrow to accidentally trigger a guess for "overcoat" or "toothpick". I will only group words if the association is extremely strong and exclusive to my team's list, avoiding any clue that could plausibly relate to the opponent's remaining words.} \\
\\ 
\texttt{Trial \#2: New plan: My previous attempt failed because I only targeted one word ("spinach") with the clue "leafy," leaving "geisha" and "cactus" unguessed, which resulted in a loss as the game requires clearing all team words. While the clue was safe, it was inefficient. Next time, I will prioritize finding a single clue that semantically links at least two of my team's words to maximize progress per turn, provided the association is strong enough to avoid distractors. For this specific set, I will attempt a clue like "green" (linking spinach and cactus) or "japan" (linking geisha and potentially spinach/cactus via cultural context if valid, though "green" is safer). If no safe multi{-}word link exists, I will immediately follow up with a second clue for the remaining word rather than stopping after one. I must ensure the TARGETS list includes every word the clue is intended to cover, and I will not stop until all three words are accounted for across my turns.} \\
\\ 
\\ 
\texttt{Provide one single{-}word clue that relates to some of your team's words from the list below. You can choose to how many and to which words your clue relates to.} \\
\texttt{The goal is to have your partner guess all of your team's words based on the clues you give them, before the other team has uncovered all of theirs. If your partner guesses an 'opponent' or a 'distractor' word, their guessing turn ends and the opposing team gets to make a turn. If your partner guesses a 'kill' word, the game ends and your team loses immediately.} \\
\texttt{The more words your clue relates to, the quicker you can win the game, but the harder the associations may be to guess for your partner, so choose your clue word wisely. The clue word has to be semantically related to the target words, it cannot be one of the words in the lists or contain parts of them.} \\
\\ 
\texttt{Always give your single{-}word clue and your comma{-}separated list of related target words in the following format and make your answers as short as possible, never include any other text than is required in this form:} \\
\\ 
\texttt{CLUE: <WORD>} \\
\texttt{TARGETS: <TARGETS>} \\
\\ 
\texttt{Your team words are: geisha, spinach, cactus.} \\
\texttt{Your opponent's team words are: .} \\
\texttt{Distractor words are: puck, cavalry, bayou, courtroom, luncheon.} \\
\texttt{The kill words are: bridegroom.} \\
            }
        }
    }
    & \\ \\

    \theutterance \stepcounter{utterance}  
    & \multicolumn{4}{p{0.6\linewidth}}{
        \cellcolor[rgb]{0.9,0.9,0.9}{
            \makecell[{{p{\linewidth}}}]{
                \texttt{\tiny{[P1$\rangle$GM]}}
                \texttt{CLUE: green} \\
\texttt{TARGETS: spinach, cactus} \\
            }
        }
    }
    & & \\ \\

    \theutterance \stepcounter{utterance}  
    & & & \multicolumn{2}{p{0.3\linewidth}}{
        \cellcolor[rgb]{0.9,0.9,0.9}{
            \makecell[{{p{\linewidth}}}]{
                \texttt{\tiny{[GM$|$GM]}}
                \texttt{green} \\
            }
        }
    }
    & & \\ \\

    \theutterance \stepcounter{utterance}  
    & & & \multicolumn{2}{p{0.3\linewidth}}{
        \cellcolor[rgb]{0.9,0.9,0.9}{
            \makecell[{{p{\linewidth}}}]{
                \texttt{\tiny{[GM$|$GM]}}
                ['spinach', 'cactus']
            }
        }
    }
    & & \\ \\

    \theutterance \stepcounter{utterance}  
    & & & \multicolumn{2}{p{0.3\linewidth}}{
        \cellcolor[rgb]{0.9,0.9,0.9}{
            \makecell[{{p{\linewidth}}}]{
                \texttt{\tiny{[GM$|$GM]}}
                {'word': 'spinach', 'assignment': 'team'}
            }
        }
    }
    & & \\ \\

    \theutterance \stepcounter{utterance}  
    & & & \multicolumn{2}{p{0.3\linewidth}}{
        \cellcolor[rgb]{0.9,0.9,0.9}{
            \makecell[{{p{\linewidth}}}]{
                \texttt{\tiny{[GM$|$GM]}}
                {'word': 'cactus', 'assignment': 'team'}
            }
        }
    }
    & & \\ \\

    \theutterance \stepcounter{utterance}  
    & & & \multicolumn{2}{p{0.3\linewidth}}{
        \cellcolor[rgb]{0.9,0.9,0.9}{
            \makecell[{{p{\linewidth}}}]{
                \texttt{\tiny{[GM$|$GM]}}
                \texttt{opponent has won} \\
            }
        }
    }
    & & \\ \\

    \theutterance \stepcounter{utterance}  
    & \multicolumn{4}{p{0.6\linewidth}}{
        \cellcolor[rgb]{0.95,0.95,0.95}{
            \makecell[{{p{\linewidth}}}]{
                \texttt{\tiny{[P2$\rightarrow$P1]}}
                \texttt{Post{-}episode Reflexion (not sent to GM):} \\
\texttt{New plan: My previous attempt failed because I provided a clue ("green") that only covered two of my three target words, leaving "geisha" unguessed and resulting in an incomplete turn. While "green" was a valid association for spinach and cactus, it was too broad and risky, potentially leading to distractor guesses like "puck" (green jersey) or "cavalry" (green uniforms), and it failed to address the full set of words. Next time, I will prioritize a clue that semantically links all three words if possible, or I will ensure my strategy explicitly accounts for covering the remaining word in the same turn or immediately following. For this specific set, I will look for a broader cultural or thematic link, such as "japan" (linking geisha and potentially cactus/spinach via Japanese cuisine or culture, though risky) or "art" (geisha as art, cactus/spinach in art). If a single safe clue for all three is impossible, I will choose a clue that covers the most words with the highest specificity to avoid distractors, and I will ensure the TARGETS list is exhaustive for that clue. Crucially, I must not stop until all team words are either guessed or covered by a sequence of clues, avoiding the mistake of leaving a word unaddressed. For this specific instance, I will try a clue like "culture" or "japan" to attempt to link all three, or if that is too ambiguous, I will use "plant" for spinach and cactus and immediately follow with a specific clue for geisha, ensuring no word is left behind. However, given the single{-}turn constraint implied by the loss, I must find a clue that covers all three or accept that I need a multi{-}step strategy where I don't declare the turn over until all are covered. The immediate fix is to ensure the TARGETS list includes all words the clue is intended to solve, and if one word is left out, I must provide a second clue or choose a different primary clue that encompasses the missing word. For "geisha, spinach, cactus", a clue like "japan" might work for geisha and spinach (Japanese food), but cactus is weak. A better approach is to realize that "green" was the best link for two, but I failed to account for the third. I will now plan to use a clue that is specific to the hardest word (geisha) first, or find a unifying theme. Actually, the most critical error was leaving "geisha" out. I will ensure my next clue attempts to cover all three, or I will structure my response to include a follow{-}up if the game format allows, but since the format is single clue per turn, I must find a clue that links all three. If no such clue exists, I must prioritize the word that is hardest to guess or most likely to be a distractor. In} \\
            }
        }
    }
    & & \\ \\

\end{supertabular}
}

\end{document}
